NB: This element has very poor parallel efficiency. Hence, the \verb|START_PASS|, \verb|END_PASS|, and
\verb|INTERVAL| options should be used to limit the frequency of computations to the minimum needed.

The output filename may be an incomplete filename.  In the case of the
\verb|SLICE| element, this means it may contain one instance of
the string format specification ``\%s'' and one occurence of an
integer format specification (e.g., ``\%ld'').  {\tt elegant} will
replace the format with the rootname (see
\verb|run_setup|) and the latter with the element's occurrence
number.  For example, suppose you had a repetitive lattice defined as
follows:
\begin{verbatim}
S1: SLICE,FILENAME=''%s-%03ld.s1''
Q1: QUAD,L=0.1,K1=1
D: DRIFT,L=1
Q2: QUAD,L=0.1,K1=-1
CELL: LINE=(S1,Q1,D,2*Q2,D,Q1)
BL: LINE=(100*CELL)
\end{verbatim}
The element \verb|S1| appears 100 times.  Each instance will result in a new file being produced.  Successive instances
have names like ``{\em rootname}-001.s1'', ``{\em rootname}-002.s1'', ``{\em rootname}-003.s1'', and so on up to ``{\em
rootname}-100.s1''.  (If instead of ``\%03ld'' you used ``\%ld'', the names would be ``{\em rootname}-1.s1'', ``{\em
rootname}-2.s1'', etc. up to ``{\em rootname}-100.s1''.  This is generally not as convenient as the names don't sort
into occurrence order.)

These files can be plotted together easily, for example,
\begin{verbatim}
sddsplot -column=Slice,ey *-???.s1 -graph=line,vary -order=spectral
\end{verbatim}

For multi-bunch simulation, a variant of this feature can be used to create multiple files, one for
each bunch. For example,
\begin{verbatim}
S1: SLICE,FILENAME=''%s-%03ld.s1'',BUNCH_SERIES=1
Q1: QUAD,L=0.1,K1=1
D: DRIFT,L=1
Q2: QUAD,L=0.1,K1=-1
CELL: LINE=(Q1,D,2*Q2,D,Q1)
BL: LINE=(100*CELL,48*S1)
\end{verbatim}
would result in the creation of 48 files, one for each of 48 bunches. If more than 48 bunches are present,
they won't be recorded. If fewer are present, the excess files will be empty.

